\chapter{$O(a)$ improvement in the Dirac Action}
% Section 13 of R. Gupta, "Introduction to Lattice QCD" (hep-lat/9807028).

The first attempt to improve the Wilson fermion action by removing the
$O(a)$ corrections was by Hamber and Wu [102].  They added a two link
term to the Wilson action
\begin{align}
\mathcal{S}_{HW}
&= A_W + \frac{r\kappa}{4a} \sum_{x,\mu} \bar{\psi}^L(x)
   \left[ U_\mu (x) U_\mu (x+\hat{\mu}) \psi^L (x+2\hat{\mu})
   + U_\mu^\dagger (x-\hat{\mu}) U_\mu^\dagger (x-2\hat{\mu}) \psi^L (x-2\hat{\mu}) \right] \nonumber\\
&\equiv \sum_{x,y} \bar{\psi}_L (x) M_{xy}^{HW} \psi_L (y) .
\label{eq:HWaction}
\end{align}

It is easy to check by Taylor expansion that, at the classical level,
the $O(a)$ artifact in $\mathcal{S}_W$ is cancelled by the HW term.  One can
include quantum corrections by calculating the relative coefficient in
perturbation theory.  It turns out that this action has not been
pursued --- in the 1980's due to the lack of computer power, and in the
1990's due to the merits of the Sheikholeslami-Wohlert (called SW or
the clover) action [20].

\section{The Sheikholeslami-Wohlert (clover) Action}

SW proposed adding a second dimension five operator to $\mathcal{S}_W$, the
magnetic moment term, to remove the $O(a)$ artifacts without
spoiling Wilson's fix for doublers.  The resulting SW action is
\begin{equation}
\mathcal{S}_{SW} \ = \ \mathcal{S}_W -
   \frac{i a C_{SW} \kappa r}{4}\ \bar{\psi}(x) \sigma_{\mu\nu} F_{\mu\nu} \psi(x) \ .
\label{eq:SWact}
\end{equation}

The fact that this action is $O(a)$ improved for on-shell quantities
cannot be deduced from a Taylor expansion.  To show improvement
requires an enumeration of all dimension 5 operators, the use of
equations of motion to remove some of them, and incorporating the
remaining two ($m^2 \bar{\psi}\psi$ and $m F^{\mu\nu} F_{\mu\nu}$)
by renormalizing $g^2 \to g^2 (1+b_g ma)$ and $m \to m (1+b_m ma)$.  The
construction of the SW action has been explained in detail by Prof.
L\"uscher at this school, so instead of repeating it, I direct you to
his lectures.  I will simply summarize results of tuning $C_{SW}$ using
both perturbative and non-perturbative methods.  A list of
perturbatively improved values of $C_{SW}$ that have been used in
simulations is as follows.

\begin{center}
\begin{tabular}{ll}
$C_{SW} = 1$                                & tree-level improvement \\
$C_{SW} = 1 + 0.2659(1) g_0^2$              & 1-loop improved [103] \\
$C_{SW} = 1/u_0^3$                          & tadpole improved \\
$C_{SW} = 1/u_0^3 (1 + 0.0159 g_I^2)$       & 1-loop tadpole improved \\
\end{tabular}
\end{center}

In the 1-loop TI result $u_0$ was chosen to be the fourth root of the
plaquette.  As an exercise convince yourself that using the expectation
value of the link in Landau gauge gives a value closer ($1.6$)
to the non-perturbatively improved result ($1.77$ at $g=1$) given
below.  For each of the above values of $C_{SW}$, the errors are reduced
from $O(a)$ (Wilson action) to at least $O(\alpha_s a)$.

The recent exciting development, also reviewed by L\"uscher, is that
the axial Ward identity $C_{SW}$ can be used to non-perturbatively
tune $C_{SW}$ and remove all $O(a)$ artifacts provided one
simultaneously improves the coupling $g^2$, the quark mass $m_q$, and
the currents [104].  Estimates of $C_{SW}$ have been made
at a number of values of $g$ for the quenched theory with Wilson's
gauge action [104,105].  These data can be fit by [105]
\begin{equation}
C_{SW} = \frac{1 - 0.6084 g^2 - 0.2015 g^4 + 0.03075 g^6}{1 - 0.8743 g^2} \ .
\label{eq:CSWnonpert}
\end{equation}
for $6/g^2 > 5.7$.

The advantages of the SW action is that it is local and leaves
perturbation theory tractable [106].  Adding the clover
term is only a $\sim 15\%$ overhead on Wilson fermion simulations.  An
example of improvement obtained using the SW versus Wilson action
for fixed quark mass ($M_\pi/M_\rho = 0.7$) is that the deviations
from the $a=0$ values in $M_N/\sqrt{\sigma}$ and $M_{\rm vector}/\sqrt{\sigma}$ at
$a=0.17$ fermi are reduced from $30-40\% \to \sim 3-5\%$
[105].  Further tests and evaluation of the improvement are
being done by a number of collaborations.  I expect a number of
quenched results to be reported at the LATTICE 98 conference, and
these should appear on the e-print archive by late 1998.

\section{D234c action}

As mentioned before, there are two sources of discretization errors:
(i) approximating derivatives by finite differences, and (ii) quantum
corrections, for example those due to modes with momentum $> \pi/a$
that are absent from the lattice theory.
Alford, Klassen, and Lepage have proposed improved gauge
and Dirac actions that reduce (i) by classical improvement and (ii)
by tadpole improvement [107,79].  Their current
proposal for isotropic lattices, called the \textbf{D234c action} as
it involves second, third, and fourth order derivatives, is
\begin{align}
M_{D234c}
&= m(1+0.5 r a m) + \sum_\mu \left\{ \gamma_\mu \Delta_\mu^{(1)}
   - \frac{C_3}{6} a^2 \gamma_\mu \Delta_\mu^{(3)} \right\} \nonumber\\
&\quad + r \sum_\mu \left\{ - \frac{1}{2} a \Delta_\mu^{(2)}
   - \frac{C_F}{4} a \sum_\nu \sigma_{\mu\nu} F_{\mu\nu}
   + \frac{C_4}{24} a^2 \Delta_\mu^{(4)} \right\} \ .
\label{eq:D234caction}
\end{align}

Here $\Delta_\mu^{(n)}$ is the $n$th order lattice covariant
derivative (symmetric and of shortest extent).  The term proportional
to $C_3$ is precisely what is needed to kill the leading
discretization error in $\Delta_\mu^{(1)}$ as shown in
Eq.~7.6.  The terms proportional to $r$ are
generated by a field redefinition and thus represent a redundant
operator.  They, therefore leave unalterated the $O(a)$ improvement of
the naive action for on-shell quantities.  The three terms in it are
respectively the Wilson term, that is again used to remove the
doublers, clover term (for $F_{\mu\nu}$ they use an $O(a^2)$ improved
version) as in the SW action needed to remove the $O(a)$ error
introduced by the Wilson fix, and an $a^3$ correction.  To include
quantum corrections to the classical values $C_3 = C_F = C_4 = 1$ they
resort to tadpole improvement, $i.e.$, they divide all links appearing
in $M_{D234c}$ by $u_0$, which is taken to be the mean link in Landau
gauge.

Quantum corrections generally also induce additional operators,
however, in perturbation theory these start at $O(\alpha_s)$.  After
tadpole improvement the coefficients of $\alpha_s$ for such operators
are all expected to be small, so one hopes that their contributions
can be neglected.  The philosophy/motivation is analogous to dropping the small
$\mathcal{L}_2^{(6)}$ term in the tadpole improved Luscher-Weisz pure gauge
action given in Eq.~12.14.

For two of the coefficients, $C_F$ and $C_3$, the authors have carried
out non-perturbative tests of the accuracy of tadpole improvement
[79].  Since only the operator $\Delta_\mu^{(3)}$ breaks
rotational symmetry at $O(a^2)$, therefore $C_3$ can be tuned
non-perturbatively by studying rotational invariance of dispersion
relations.  Similary, the clover term and the Wilson term break chiral
symmetry at $O(a)$, thus their relative coefficient $C_F$ can be tuned
by studying the axial Ward identity.  Results of their tests show that
errors in the spectrum and dispersion relations for momenta as high as
$p = 1.5/a$ are $\approx 10\%$ even at $a = 0.4$ fermi!  The authors
are extending their analyses of the $D234$ action to anisotropic
lattices which, as will be discussed in Section 18.4,
have already yielded rich dividends in the calculation of glueball
masses.

\section{NRQCD action}

In the simulation of heavy quarks one has to worry about $O(m_H a)$
errors in addition to $O(pa)$ and $O(\Lambda_{\rm QCD} a)$.  For charm and
bottom quarks these errors are $O(1)$ in the range of lattice spacings
one can reasonably simulate.  Lepage and collaborators, recognizing
that the heavy quark in heavy-heavy or heavy-light systems is, to a
good approximation, non-relativistic, developed a non-relativistic
formulation (NRQCD) valid for $m_H a > 1$ [108].
In this approach, the action for the two-component Pauli spinors
is obtained by a Foldy-Wouthuysen transformation, and relativistic and
discretization corrections are added as a systematic expansion in $1/M$
and $p/M$.
One version of the NRQCD action is [109]
\begin{align}
\mathcal{L}
&= \bar{\psi}_t \psi_t \nonumber\\
&\quad - \bar{\psi}_t
   \left( 1 - \frac{a \delta H}{2} \right)_t
   \left( 1 - \frac{a H_0}{2n} \right)^n_t
   U_4^\dagger
   \left( 1 - \frac{a H_0}{2n} \right)^n_{t-1}
   \left( 1 - \frac{a \delta H}{2} \right)_{t-1} \psi_{t-1},
\label{eq:nrqcdact}
\end{align}
where $\psi$ is the two-component Pauli spinor and
$H_0$ is the nonrelativistic kinetic energy operator,
\begin{equation}
H_0 = - \frac{\Delta^{(2)}}{2 M_0} \ .
\label{eq:nrqcdH0}
\end{equation}
$\delta H$ includes relativistic and finite-lattice-spacing
corrections,
\begin{align}
\delta H
&= - \frac{g c_1}{2 M_0} \boldsymbol{\sigma}\cdot\boldsymbol{B} \nonumber\\
&\quad + \frac{i g c_2}{8 (M_0)^2}
   \left( \boldsymbol{\Delta}\cdot\boldsymbol{E} - \boldsymbol{E}\cdot\boldsymbol{\Delta} \right)
   - \frac{g c_3}{8 (M_0)^2}
   \boldsymbol{\sigma}\cdot
   \left( \boldsymbol{\Delta}\times\boldsymbol{E} - \boldsymbol{E}\times\boldsymbol{\Delta} \right) \nonumber\\
&\quad - \frac{c_4 (\Delta^{(2)})^2}{8 (M_0)^3}
   + \frac{c_5 a^2 \Delta^{(4)}}{24 M_0}
   - \frac{c_6 a (\Delta^{(2)})^2}{16 n (M_0)^2} \ .
\label{eq:NRQCDdeltaH}
\end{align}
Here $\boldsymbol{E}$ and $\boldsymbol{B}$ are the electric and magnetic fields,
$\boldsymbol{\Delta}$ and $\Delta^{(2)}$ are the
spatial gauge-covariant lattice derivative and
laplacian, and $\Delta^{(4)}$ is a lattice version of the continuum
operator $\sum_i D_i^4$.  The parameter $n$ is introduced to remove
instabilities in the heavy quark propagator due to the highest
momentum modes [108].  Again, the tree-level
coefficients $c_i = 1$ can be tadpole-improved by rescaling all the
link matrices by $u_0$.

This formulation has been used very successfully to calculate
$\alpha_s$ from the upsilon spectrum [110] as discussed
in Section 19, the heavy-heavy and heavy-light spectrum
[111], and the heavy-light decay constants $f_B$ and $f_{B_s}$
[109].  The drawback of NRQCD is that it is an effective theory
and one cannot take $a \to 0$.  In fact, to simulate charm and bottom
quarks, one has to work in a narrow window $1 \lesssim 1/a \lesssim 2.5$
GeV to keep $m_H a > 1$.  The good news is that current data
[110,111,109] show that
the discretization and relativistic corrections are small and under
control with the terms included in Eq.~13.7.

\section{Perfect action}

The development of a perfect Dirac action is not complete.  The method
is same as used for the gauge action and discussed in [90].
The bottleneck is that the saddle point integration gives an action that even
in the free field limit involves a number of sites, and the gauge
connections between any pair of sites is the average of a large number
of paths [112].  First tests of such an action truncated to
a $3^4$ hypercube have been carried out by Degrand and collaborators
[100].  From these it is not clear whether the
improvement, compared to other discretization schemes, justifies the
extra cost of simulating such an action.  Hopefully better ways to tune
and test such actions will be formulated soon.

\section{Fermilab action}

The Fermilab group has proposed a ``Heavy Wilson'' fermion action
[113].  This action is meant to interpolate between the
light $O(a)$ improved SW-clover action and the NRQCD action for heavy
quarks.  The goal is to remove not only $O(a)$ discretization
corrections but also take into account all powers of $ma$ errors so
that charmonium and bottomonium systems can be simulated with the same
reliability as light-light physics.  This improvement is especially
targeted towards improving heavy-light matrix elements.  So far no
tests have been reported with the proposed action, therefore I will
not reproduce the action here but refer interested readers to
[113] for details.

\section{Summary}

In analogy to discussions of improved gauge actions, having seen
tangible evidence for improvement in quenched lattice data with each of
the above formulations, the question is, which action is the most
efficient for a comprehensive study?  The deciding factor, in my
opinion, will be dynamical simulations since the prefactor and the
scaling exponent ($L^{\sim 10}$) makes the cost of update of
background configurations prohibitive without improvement on today's
computers.  Dynamical simulations using improved actions have just
begun, so it will be a few years before enough data can be generated
to determine the advantages and disadvantages of different approaches.

Finally, I should mention a recent novel approach --- domain wall
fermions --- for improving the chiral properties of Wilson fermions
[114].  In this approach the system (domain) is
5-dimensional and the fermions live on 4-dimensional surfaces that
make up the opposite boundaries (walls).  The 4-dimensional gauge
fields on the two walls are identical, whereas the fermion zero-modes
have opposite chirality.  The overlap between these modes is
exponentially suppressed by the separation in the fifth
dimension.  Consequently, the chiral symmetry breaking contributions are
also supressed, and results with good chiral properties for vector
like theories like QCD can be obtained by adding the contribution from
the two domain walls.  First simulations of weak matrix elements show
results with good chiral properties and in general look extremely
promising [115].  Since I do
not do justice to this exciting approach in these lectures, interested
readers should consult [115] and references therein.
